% Compile from slide/20261007/ (PDF: slide/out/20261007/ , other files: build/):
%   ../lualatex-dated.sh -synctex=1 -interaction=nonstopmode -file-line-error -output-directory=out survey_kawamotostyle_20261007_handout_v3_en.tex
\documentclass[aspectratio=169,10pt,t,handout]{beamer}
\usepackage[en]{labslide}
% Raise every requested font size by 1pt, as in the Japanese handout.
\makeatletter
\newdimen\labvtwo@size
\newskip\labvtwo@skip
\let\labvtwo@set@fontsize\set@fontsize
\def\set@fontsize#1#2#3{%
  \@defaultunits\labvtwo@size#2pt\relax\@nnil
  \advance\labvtwo@size by 1pt\relax
  \@defaultunits\labvtwo@skip#3pt\relax\@nnil
  \advance\labvtwo@skip by 1pt\relax
  \begingroup
    \edef\x{\endgroup
      \noexpand\labvtwo@set@fontsize{\unexpanded{#1}}{\strip@pt\labvtwo@size}{\the\labvtwo@skip}%
    }\x
}
\makeatother
\usetikzlibrary{fit,backgrounds,calc}

% Place the title midway between the top of the page and the rule below it.
\makeatletter
\setlength{\beamer@frametopskip}{0mm}
\makeatother
\setbeamertemplate{frametitle}{%
  \vspace{4.2mm}%
  {\usebeamerfont{frametitle}\color{ink}\insertframetitle\par}%
  \vspace{0.4mm}%
  {\color{ruleblue}\rule{\textwidth}{0.7pt}\par}%
  \vspace{1.1mm}%
  \ifx\insertframesubtitle\@empty\else
    {\footnotesize\color{muted}\insertframesubtitle\par}%
    \vspace{1.2mm}%
  \fi
}

\definecolor{cviolet}{HTML}{F3E8F6}
\definecolor{cviolettt}{HTML}{7A458E}
% Same pastel fill and mid-tone title as the other cards. Red is not an alias of orange.
\definecolor{cred}{HTML}{F8E6E6}
\definecolor{credtt}{HTML}{C44545}
\renewcommand{\usecardcolor}[1]{%
  \def\cardfill{cgreen}%
  \def\cardtitlec{cgreentt}%
  \ifstrequal{#1}{red}{%
    \def\cardfill{cred}%
    \def\cardtitlec{credtt}%
  }{}%
  \ifstrequal{#1}{orange}{%
    \def\cardfill{corange}%
    \def\cardtitlec{corangett}%
  }{}%
  \ifstrequal{#1}{blue}{%
    \def\cardfill{cblue}%
    \def\cardtitlec{cbluett}%
  }{}%
  \ifstrequal{#1}{green}{%
    \def\cardfill{cgreen}%
    \def\cardtitlec{cgreentt}%
  }{}%
  \ifstrequal{#1}{violet}{%
    \def\cardfill{cviolet}%
    \def\cardtitlec{cviolettt}%
  }{}%
  \ifstrequal{#1}{plain}{%
    \def\cardfill{cplain}%
    \def\cardtitlec{cplaintt}%
  }{}%
  \ifstrequal{#1}{gray}{%
    \def\cardfill{cplain}%
    \def\cardtitlec{cplaintt}%
  }{}%
}

\graphicspath{{fig/survey/}}

\kicker{Survey}
\title{A local planner that chooses\\easy ground for a quadruped}
\date{October 2026}

% Bullets inside cards. The second optional argument matches heights in a row.
\newcommand{\li}{%
  \ifvmode\else\par\fi
  \noindent\hangindent=1.15em\hangafter=1
  \makebox[1.15em][l]{\textcolor{ruleblue}{\textbullet}}\ignorespaces}
\newcommand{\lii}{%
  \ifvmode\else\par\fi
  \noindent\hspace{1.0em}\hangindent=2.15em\hangafter=1
  \makebox[1.05em][l]{\textcolor{muted}{\textendash}}\ignorespaces}

% Cancel the \topskip gap above the first card, then share the remaining space.
\newcommand{\cardbalance}{\par\vspace*{-\topskip}\vspace{0pt plus 1fill}}
\newcommand{\takeawaybalance}[1]{%
  \par\vspace{0.15em plus 1fill}%
  \begingroup
  \setlength{\fboxsep}{5pt}%
  \noindent\colorbox{ink}{%
    \parbox[c]{\dimexpr\textwidth-10pt\relax}{%
      \color{white}\footnotesize #1}}%
  \endgroup
}

\RenewDocumentCommand{\infocard}{O{green} o m m}{%
  \usecardcolor{#1}%
  \IfValueTF{#2}{\def\cardminh{#2}}{\def\cardminh{0pt}}%
  \begin{tikzpicture}[baseline=(card.north)]
    \node[
      anchor=north west,
      fill=\cardfill,
      text width=\dimexpr\linewidth-14pt\relax,
      minimum height=\cardminh,
      inner xsep=6.5pt,
      inner ysep=4pt,
      align=flush left
    ] (card) {%
      {\color{\cardtitlec}\bfseries #3}\par\vspace{0.15em}%
      {\color{body}\footnotesize #4\par}%
    };
  \end{tikzpicture}%
}

\begin{document}

\begin{frame}[plain]
  \labtitlepage{
    B4, Hiromu Saito, 2026/10/07
  }
\end{frame}

\begin{frame}{Research theme | Robot control}
  \centering
  \begin{tikzpicture}[node distance=6.5mm]
    \node[kworange, text width=28mm, minimum height=8mm] (c) {Controller};
    \node[kwgreen, text width=28mm, minimum height=8mm, right=of c] (p) {Local planner};
    \node[kwblue, text width=28mm, minimum height=8mm, right=of p] (g) {Gait model};
    \node[kwpurple, text width=24mm, minimum height=8mm, right=of g] (m) {Motor};
    \draw[kwarr] (c) -- (p)
      node[midway, below=4mm, font=\scriptsize, text=muted] {target pose};
    \draw[kwarr] (p) -- (g)
      node[midway, below=4mm, font=\scriptsize, text=muted] {target speed};
    \draw[kwarr] (g) -- (m)
      node[midway, below=4mm, font=\scriptsize, text=muted] {torque};
  \end{tikzpicture}

  \vspace{1.5em}
  \includegraphics[width=\textwidth,height=35mm,keepaspectratio]{img001.png}

  \takeaway{In general, the input reaches the motors through three transforms.}
\end{frame}

\begin{frame}{Research theme | Local planner}
  \centering
  \begin{tikzpicture}[node distance=6.5mm]
    \node[kworange, text width=28mm, minimum height=8mm, opacity=0.2] (c) {Controller};
    \node[kwgreen, text width=28mm, minimum height=8mm, right=of c] (p) {Local planner};
    \node[kwblue, text width=28mm, minimum height=8mm, right=of p] (g) {Gait model};
    \node[kwpurple, text width=24mm, minimum height=8mm, right=of g, opacity=0.2] (m) {Motor};
    \draw[kwarr] (c) -- (p)
      node[midway, below=4mm, font=\scriptsize, text=muted] {target pose};
    \draw[kwarr] (p) -- (g)
      node[midway, below=4mm, font=\scriptsize, text=muted] {target speed};
    \draw[kwarr] (g) -- (m)
      node[midway, below=4mm, font=\scriptsize, text=muted] {torque};
  \end{tikzpicture}

  % img2
  \includegraphics[width=\textwidth,height=35mm,keepaspectratio]{img002.png}

  \takeaway{The target speed should cover both safety and preference.}
\end{frame}

\begin{frame}{Papers on preference}
  \infocard[green][22mm]{Navigation from a walking-cost map}{%
    \li Where Should I Walk? Predicting Terrain Properties From Images Via Self-Supervised Learning
    \li IEEE RA-L, 2019
    \li Cited by 268}
  \vspace{0.4em}
  \infocard[blue][22mm]{Navigation from RGB-D and reinforcement learning}{%
    \li Traversability-Aware Legged Navigation by Learning from Real-World Visual Data
    \li IEEE Transactions on Robotics
    \li Cited by 13}

  \takeaway{The first paper puts preference in a map. The second puts it in navigation.}
\end{frame}

% ---------- Wellhausen ----------

\begin{frame}{Where Should I Walk?}
  \source{Wellhausen et al., IEEE RA-L, 2019. Cited by 268}
  \centering
  \includegraphics[width=\textwidth,height=52mm,keepaspectratio]{wsiw_row.png}

  \takeaway{RGB in, a cost map out, and Dijkstra draws the path.}
\end{frame}

\begin{frame}{Background and aim}
  \source{Wellhausen et al., IEEE RA-L, 2019}
  \cardbalance
  \infocard[orange]{Background}{%
    \li Natural ground is not only rigid
    \li Even when flat, water, sand, and soil change with moisture
    \lii A class defined by people still varies in its properties
    \lii Manual annotation is very hard, or impossible}
  \vspace{0.7em plus 1fill}
  \infocard[green]{Aim}{%
    \li Tie an RGB image to what the foot feels
    \li Predict a walking cost from the RGB image}
  \takeawaybalance{Tie foot sensation to the image, and predict a walking cost at each pixel.}
\end{frame}

\begin{frame}{Contributions}
  \source{Wellhausen et al., IEEE RA-L, 2019}
  \cardbalance
  \infocard[blue]{Contribution 1}{\li Tie the camera image to the foot's contact force}
  \vspace{0.35em}
  \infocard[orange]{Contribution 2}{\li Predict terrain properties while walking. The metric is the ground reaction score}
  \vspace{0.35em}
  \infocard[green]{Contribution 3}{\li A CNN trained on sparse labels}
  \vspace{0.35em}
  \infocard[violet]{Contribution 4}{\li Intuitive navigation was achieved}
  \takeawaybalance{Foot-sensor features become a ground reaction score, and that cost chooses the path.}
\end{frame}

\begin{frame}{Proposed method}
  \source{Wellhausen et al., IEEE RA-L, 2019}
  \centering
  \begin{tikzpicture}[
      scale=0.82, transform shape,
      font=\footnotesize,
      box/.style={kwblue, text width=32mm, minimum height=12mm, inner sep=2.5pt},
      tag/.style={font=\scriptsize\bfseries, text=muted}
    ]
    \node[box] (lab) at (0,0) {Foothold labeling};
    \node[box] (weak) at (4.35,1.15) {Weak labels\\{\scriptsize (Segmentation)}};
    \node[box] (self) at (4.35,-1.15) {Self-supervised labels\\{\scriptsize (Ground Reaction Score)}};
    \node[box] (cls) at (8.7,1.15) {Terrain-class estimator};
    \node[box] (cost) at (8.7,-1.15) {Walking-cost map};
    \node[box] (corr) at (13.05,1.15) {Check the correlation};
    \node[box] (nav) at (13.05,-1.15) {Navigation\\{\scriptsize (Dijkstra)}};
    \begin{scope}[on background layer]
      \node[fill=cblue, rounded corners=1.6pt, inner sep=2.2mm, fit=(lab)] (p1) {};
      \node[fill=corange, rounded corners=1.6pt, inner sep=2.2mm, fit=(weak)(self)] (p2) {};
      \node[fill=cgreen, rounded corners=1.6pt, inner sep=2.6mm, fit=(cls)(cost)] (p3) {};
      \node[fill=cviolet, rounded corners=1.6pt, inner sep=2.2mm, fit=(nav)] (p4) {};
    \end{scope}
    \node[tag, text=cbluett, anchor=south] at (p1.north) {Contrib.\ 1};
    \node[tag, text=corangett, anchor=south] at (p2.north) {Contrib.\ 2};
    \node[tag, text=cgreentt, anchor=south] at (p3.north) {Contrib.\ 3};
    \node[tag] at ($(p3.north)+(0,-20.8mm)$) {Mean Teacher};
    \node[tag, text=cviolettt, anchor=south] at (p4.north) {Contrib.\ 4};
    \draw[kwarr] (lab.east) -- (weak.west);
    \draw[kwarr] (lab.east) -- (self.west);
    \draw[kwarr] (weak) -- (cls);
    \draw[kwarr] (self) -- (cost);
    \draw[kwarr] (cls) -- (corr);
    \draw[kwarr] (cost) -- (nav);
  \end{tikzpicture}

  \takeaway{Weak labels check the terrain class. The self-supervised cost map chooses the path.}
\end{frame}

\begin{frame}{Contribution 1 | Foothold labeling}
  \source{Wellhausen et al., IEEE RA-L, 2019}
  \centering
  {\footnotesize Transform the contact center and two points on the ellipse, three points in all, onto the image.\par}
  \vspace{0.15em}
  {\footnotesize
  \[
    \bigl[\, \mathbf{x}^{i,t}_{c} \quad \mathbf{x}^{i,t}_{1} \quad \mathbf{x}^{i,t}_{2} \,\bigr]
    = \mathbf{K} \cdot \mathbf{T}^{t}_{CW} \cdot
    \bigl[\, \mathbf{p}^{i}_{c} \quad \mathbf{p}^{i}_{1} \quad \mathbf{p}^{i}_{2} \,\bigr].
  \]
  }
  \vspace{-0.55em}
  {\footnotesize $K$: camera intrinsics. The matrix for the 3D-to-2D transform.\par}

  \vspace{0.45em}
  \includegraphics[width=\textwidth,height=32mm,keepaspectratio]{img003.png}

  \takeaway{The transform and the RGB-D pose history locate the foothold in each image.}
\end{frame}

\begin{frame}{Labeling for weak supervision}
  \source{Wellhausen et al., IEEE RA-L, 2019}
  \begin{columns}[T,onlytextwidth]
    \begin{column}{0.50\textwidth}
      \infocard[orange][32mm]{A person watches the video}{%
        \li Label the intervals where the robot walks
        \lii Asphalt, Path, Grass, Dirt, Sand
        \li That frame's foothold receives the label
        \li Pixels in the image receive the same label}
    \end{column}
    \hfill
    \begin{column}{0.50\textwidth}
      \centering
      \includegraphics[width=\linewidth,height=45mm,keepaspectratio]{class_dots.png}\\[0.15em]
      {\scriptsize\color{muted} Grass, Path, Sand}
    \end{column}
  \end{columns}

  \takeaway{Weak supervision is for visualization and for evaluation.}
\end{frame}

\begin{frame}{Contribution 2 | Self-supervised ground reaction score}
  \source{Wellhausen et al., IEEE RA-L, 2019}
  \begin{columns}[T,onlytextwidth]
    \begin{column}{0.50\textwidth}
      {\footnotesize\color{body} Define the Ground Reaction Score (GRS) at a point.}
      \vspace{0.35em}
      \infocard[blue][13mm]{1}{\li Torque-sensor series during contact}
      \vspace{0.28em}
      \infocard[orange][13mm]{2}{\li A wavelet transform yields 13860 dimensions}
      \vspace{0.28em}
      \infocard[green][13mm]{3}{\li The first principal component of PCA is the GRS}
    \end{column}
    \hfill
    \begin{column}{0.50\textwidth}
      \centering
      \includegraphics[width=\linewidth,height=48mm,keepaspectratio]{grs_dots.png}
    \end{column}
  \end{columns}

  \takeaway{The image shows that the GRS differs by material.}
\end{frame}

\begin{frame}{Contribution 3 | Labels are sparse, so use Mean Teacher}
  \source{Wellhausen et al., IEEE RA-L, 2019}
  \vspace*{-0.4mm}
  \begin{columns}[T,onlytextwidth]
    \begin{column}{0.50\textwidth}
      \infocard[orange][22mm]{Too sparse to train}{%
        \li Most labels are at the feet. The upper image has little data
        \li At most 10\% of the pixels are labeled}
      \vspace{0.18em}
      \infocard[blue][22mm]{Mean Teacher}{%
        \li Different augmentations of the same image
        \li Each augmentation goes into a different model
        \li The image is the same, so the two outputs should match}
    \end{column}
    \hfill
    \begin{column}{0.50\textwidth}
      \centering
      \includegraphics[width=\linewidth,height=48mm,keepaspectratio]{sparse.png}
    \end{column}
  \end{columns}

  \takeaway{Foothold points alone are sparse, so training matches the two augmented outputs.}
\end{frame}

\begin{frame}{Mean Teacher}
  \source{Wellhausen et al., IEEE RA-L, 2019}
  \centering
  \begin{tikzpicture}[scale=0.74, transform shape, box/.style={minimum height=10.5mm, text width=28mm, font=\normalsize}]
    \node[kw, draw=cplaintt, fill=cplain, box] (img) at (0,0) {Image};
    \node[kworange, box] (s) at (5.5,1.48) {Student Model};
    \node[kwgreen, box] (t) at (5.5,-1.48) {Teacher Model};
    \node[kwblue, box] (loss) at (10.2,0) {Difference\\[0.1em] LOSS};
    \draw[kwarr] (img.east) -- (s.west)
      node[pos=0.42, above=1.6mm, font=\footnotesize, text=muted] {Augmentation A};
    \draw[kwarr] (img.east) -- (t.west)
      node[pos=0.42, below=1.6mm, font=\footnotesize, text=muted] {Augmentation B};
    \draw[kwarr] (s.east) -- (loss.north west);
    \draw[kwarr] (t.east) -- (loss.south west);
    \draw[kwarr, cgreentt] (s.south) -- (t.north)
      node[midway, right=1.6mm, font=\footnotesize, text=cgreentt] {EMA};
  \end{tikzpicture}

  \vspace{0.08em}
  \begin{columns}[T,onlytextwidth]
    \begin{column}{0.32\textwidth}
      \infocard[blue][14mm]{Augmentation}{\li Augment the same image in two ways \li Feed the two models}
    \end{column}
    \hfill
    \begin{column}{0.32\textwidth}
      \infocard[orange][14mm]{Student}{\li The image is the same, so the output should match the Teacher. \li The difference is the loss}
    \end{column}
    \hfill
    \begin{column}{0.32\textwidth}
      \infocard[green][14mm]{Teacher}{\li No direct gradient step. \li Update by a moving average of the Student weights}
    \end{column}
  \end{columns}

  \takeaway{Augmentation gives a target to unlabeled pixels, so they can be trained.}
\end{frame}

\begin{frame}{Results | Weakly supervised classification}
  \source{Wellhausen et al., IEEE RA-L, 2019, Table~I}
  \vspace{0.2em}
  \centering
  {\footnotesize
  \begin{tabular}{lcccccc}
    \toprule
    & asphalt & path & grass & dirt & sand & Mean \\
    \midrule
    Accuracy & 93.1\% & 97.0\% & 97.3\% & 85.2\% & 80.3\% & \textbf{90.6\%} \\
    \bottomrule
  \end{tabular}
  }

  \vspace{0.7em}
  \infocard[green]{Overall}{\li Mean accuracy is 90.6\%.}
  \vspace{0.4em}
  \infocard[orange]{dirt and sand}{\li Few samples, which may be why the accuracy is lower}

  \takeaway{For labels this sparse, the result is very strong.}
\end{frame}

\begin{frame}{Results | Self-supervised regression}
  \source{Wellhausen et al., IEEE RA-L, 2019}
  \begin{columns}[T,onlytextwidth]
    \begin{column}{0.45\textwidth}
      \centering
      \includegraphics[width=\linewidth,height=42mm,keepaspectratio]{pca.png}
    \end{column}
    \hfill
    \begin{column}{0.55\textwidth}
      \infocard[green][22mm]{Separated by material}{%
        \li The first principal component differs by terrain material}
      \vspace{0.4em}
      \infocard[orange][22mm]{Stiffer ground varies more}{%
        \li Stiffer ground vibrates more at the instant of contact
        \li The vibration may not be fully absorbed}
    \end{column}
  \end{columns}

  \takeaway{The first component separates by material, and stiffer materials vary more.}
\end{frame}

\begin{frame}{Results | Class and ground reaction score}
  \source{Wellhausen et al., IEEE RA-L, 2019}
  \centering
  \includegraphics[height=55mm]{img004.png}\hspace{1.2em}%
  \includegraphics[height=20mm]{img005.png}

  \takeaway{Terrain is classified by material, and each material has a characteristic predicted GRS.}
\end{frame}

\begin{frame}{Contribution 4 | Navigator}
  \source{Wellhausen et al., IEEE RA-L, 2019}
  \begin{columns}[T,onlytextwidth]
    \begin{column}{0.50\textwidth}
      \vspace{2em}
      \infocard[green][30mm]{Experiment}{%
        \li 6\,Hz on a Jetson TX2
        \li Predict the ground reaction score on a 10\,cm 2D grid
        \li Always command the robot to go 10\,m ahead
        \li Compute the path by Dijkstra}
    \end{column}
    \hfill
    \begin{column}{0.50\textwidth}
      \centering
      \includegraphics[width=\linewidth,height=52mm,keepaspectratio]{navpath.png}
    \end{column}
  \end{columns}

  \takeaway{It chose an easy path and walked along it.}
\end{frame}

\begin{frame}{Summary | Where Should I Walk?}
  \source{Wellhausen et al., IEEE RA-L, 2019}
  \cardbalance
  \begin{columns}[T,onlytextwidth]
    \begin{column}{0.49\textwidth}
      \infocard[green][52mm]{Contributions}{%
        \li Tie the camera image to the foot's contact force
        \li Predict terrain properties while walking
        \lii Ground reaction score
        \li A CNN trained on sparse labels
        \li Intuitive navigation}
    \end{column}
    \hfill
    \begin{column}{0.49\textwidth}
      \infocard[orange][25mm]{Limitations}{%
        \li Weak in dynamic environments
        \lii It looks only a short distance ahead
        \lii Not suited to moving obstacles}
      \vspace{0.4em}
      \infocard[blue][25mm]{Outlook}{%
        \li Put the walking cost into the gait controller
        \lii This paper put it into path planning
        \lii It might also set step length and step height}
    \end{column}
  \end{columns}
\end{frame}

\begin{frame}{Papers on preference}
  \infocard[green][22mm]{Navigation from a walking-cost map}{%
    \li Where Should I Walk? Predicting Terrain Properties From Images Via Self-Supervised Learning
    \li IEEE RA-L, 2019
    \li Cited by 268}
  \vspace{0.4em}
  \infocard[blue][22mm]{Navigation from RGB-D and reinforcement learning}{%
    \li Traversability-Aware Legged Navigation by Learning from Real-World Visual Data
    \li IEEE Transactions on Robotics
    \li Cited by 13}

  \takeaway{The first paper puts preference in a map. The second puts it in navigation.}
\end{frame}

% ---------- Zhang et al. ----------

\begin{frame}{Traversability-Aware Legged Navigation}
  \source{Zhang et al., IEEE Transactions on Robotics. Cited by 13}

  \vspace{0.25em}
  \centering
  \includegraphics[width=\textwidth,height=48mm,keepaspectratio]{trav_teaser.png}

  \takeaway{An RL local planner that outputs velocity directly from RGB and depth.}
\end{frame}

\begin{frame}{Background and aim}
  \source{Zhang et al., IEEE Transactions on Robotics}
  \infocard[orange][24mm]{Background}{%
    \li Obstacles can be detected, and whether they can be crossed can be judged
    \li Friction and compliance are hard to predict from appearance
    \li Physical properties are known only by walking the terrain}
  \vspace{0.35em}
  \infocard[green][18mm]{Aim}{%
    \li Autonomous quadruped navigation outdoors that accounts for traversability, by reinforcement learning}

  \takeaway{From walking experience, choose actions that stay on easy ground.}
\end{frame}

\begin{frame}{Contributions}
  \source{Zhang et al., IEEE Transactions on Robotics}
  \infocard[blue][13mm]{Contribution 1}{%
    \li Traversability from self-experience, with uncertainty}
  \vspace{0.28em}
  \infocard[green][13mm]{Contribution 2}{%
    \li Navigation that avoids terrain, using RGB-D together}
  \vspace{0.28em}
  \infocard[orange][13mm]{Contribution 3}{%
    \li Sample-efficient learning on the real robot, by an online and offline pipeline}

  \takeaway{Traversability is evaluated from self-experience and uncertainty.}
\end{frame}

\begin{frame}{Problem setting}
  \source{Zhang et al., IEEE Transactions on Robotics}
  \infocard[blue][12mm]{1. Reach the goal}{\li Arrive in minimum time}
  \vspace{0.28em}
  \infocard[orange][12mm]{2. Avoid obstacles}{\li Avoid dynamic and static obstacles}
  \vspace{0.28em}
  \infocard[green][12mm]{3. Traversal cost}{\li From its own experience, choose easy ground}

  \takeaway{The aim is to arrive, to avoid, and to stay on easy ground, all at once.}
\end{frame}

\begin{frame}{Proposed method | Architecture}
  \source{Zhang et al., IEEE Transactions on Robotics}
  \vspace*{-\topskip}
  \infocard[blue][16mm]{Architecture}{%
    \li A local planner $\pi^{p}_{\mathrm{rgbd}}$ that takes RGB-D and the goal
    \li A blind locomotion planner $\pi^{c}$
    \li A Traversability Estimator that predicts the traversability cost used as the training reward}

  \vspace{0.2em}
  \centering
  \begin{tikzpicture}[
      scale=0.78, transform shape,
      font=\scriptsize,
      inp/.style={kwblue, minimum width=20mm, minimum height=7.2mm, inner sep=1.6pt},
      mid/.style={kwblue, align=center, inner sep=2pt}
    ]
    \node[inp] (rgb) at (0,1.05) {RGB};
    \node[inp] (dep) at (0,0) {Depth};
    \node[inp] (goal) at (0,-1.05) {Goal};
    \node[mid, minimum width=24mm, minimum height=26mm, text width=22mm] (p) at (3.15,0)
      {Local planner\\$\pi^{p}_{\mathrm{rgbd}}$};
    \node[mid, minimum width=24mm, minimum height=9mm, text width=22mm] (walk) at (6.55,0.62)
      {Locomotion\\$\pi^{c}$};
    \node[kwgreen, align=center, minimum width=26mm, minimum height=9mm, text width=24mm, inner sep=2pt] (est) at (6.55,-0.72)
      {Traversability\\Estimator};
    \node[mid, minimum width=18mm, minimum height=26mm, text width=16mm] (rob) at (9.55,0) {Robot};
    \node[kwgreen, align=center, minimum width=46mm, minimum height=7mm, inner sep=2pt] (buf) at (4.7,-2.45)
      {Replay buffer: $o^{p},a^{p},r^{p}$};
    \draw[kwarr] (rgb.east) -- (rgb.east -| p.west);
    \draw[kwarr] (dep.east) -- (p.west);
    \draw[kwarr] (goal.east) -- (goal.east -| p.west);
    \draw[kwarr] (p.east |- walk.west) -- (walk.west);
    \draw[kwarr] (p.east |- est.west) -- (est.west);
    \draw[kwarr] (walk.east) -- (walk.east -| rob.west);
    \draw[kwarr] (rob.west |- est.east) -- (est.east);
    \coordinate (bus) at ($(buf.north)+(0,3.2mm)$);
    \draw[cgreentt] (goal.south) -- (goal.south |- bus);
    \draw[cgreentt] (p.south) -- (p.south |- bus);
    \draw[cgreentt] (est.south) -- (est.south |- bus);
    \draw[cgreentt] (goal.south |- bus) -- (est.south |- bus);
    \draw[kwarr, cgreentt] (buf.north |- bus) -- (buf.north);
  \end{tikzpicture}

  \takeaway{The training reward includes the traversability cost.}
\end{frame}

\begin{frame}{Proposed method | Four stages}
  \source{Zhang et al., IEEE Transactions on Robotics}
  \centering
  \includegraphics[width=\textwidth,height=52mm,keepaspectratio]{stages.png}

  \takeaway{Order: locomotion, Traversability Estimator, depth planner, real-world collection, online learning.}
\end{frame}

\begin{frame}{Stage 1 | Locomotion and traversability $T$ (Contribution 1)}
  \source{Zhang et al., IEEE Transactions on Robotics}
  \begin{columns}[T,onlytextwidth]
    \begin{column}{0.44\textwidth}
      \infocard[blue]{Locomotion $\pi^{c}$}{%
        \li Train with PPO to track the target velocity}
      \vspace{0.35em}
      \infocard[green]{Estimator $T$}{%
        \li Use the critic value $V$
        \lii Easier ground has higher value
        \li The cost is the gap between flat ground and the current terrain
        \lii This ignores differences in the command
        \li When the variance of $V_{\mathrm{terrain}}$ is large, use the tracking error $T_{\mathrm{track}}$ as the cost}
    \end{column}
    \hfill
    \begin{column}{0.54\textwidth}
      \centering
      \includegraphics[width=\linewidth,height=26mm,keepaspectratio]{stage1.png}\\[0.35em]
      \includegraphics[width=\linewidth,height=24mm,keepaspectratio]{eq1.png}
    \end{column}
  \end{columns}

  \takeaway{Easier ground, closer to flat, has higher value. Large variance falls back to the tracking cost.}
\end{frame}

\begin{frame}{Stage 2 | Depth-only planner}
  \source{Zhang et al., IEEE Transactions on Robotics}
  \vspace*{-0.8\baselineskip}
  \begin{columns}[T,onlytextwidth]
    \begin{column}{0.42\textwidth}
      \infocard[plain]{Input}{%
        \li Depth
        \li Goal}
      \vspace{0.32em}
      \infocard[blue]{Output}{\li Change in velocity}
      \vspace{0.32em}
      \infocard[green]{Training}{%
        \li Train in simulation
        \li 2a: inner product of the goal direction and the velocity
        \li 2b: forward speed, because the rear is not visible
        \li 2c: collision penalty}
    \end{column}
    \hfill
    \begin{column}{0.55\textwidth}
      \centering
      \includegraphics[width=\linewidth,height=22mm,keepaspectratio]{stage2.png}\\[0.3em]
      \includegraphics[width=\linewidth,height=20mm,keepaspectratio]{eq2.png}
    \end{column}
  \end{columns}

  \takeaway{Before the real world, a depth planner learns the goal and the obstacles from depth alone.}
\end{frame}

\begin{frame}{Stage 3 | Collect real-world data}
  \source{Zhang et al., IEEE Transactions on Robotics}
  \begin{columns}[T,onlytextwidth]
    \begin{column}{0.46\textwidth}
      \infocard[orange][18mm]{Method}{%
        \li Walk with the Stage 2 depth planner
        \li A person commands the velocity}
      \vspace{0.12em}
      \infocard[green][30mm]{Data collected}{%
        \li State $o^{p}_{\mathrm{rgbd}}$
        \lii RGB features, depth features
        \lii Robot state, command toward the goal
        \li Action $a^{p}$: change in velocity
        \li Reward: a reward that includes the traversability cost}
    \end{column}
    \hfill
    \begin{column}{0.51\textwidth}
      \centering
      \includegraphics[width=\linewidth,height=32mm,keepaspectratio]{stage3.png}\\[0.1em]
      \includegraphics[width=\linewidth,height=12mm,keepaspectratio]{eq3.png}
    \end{column}
  \end{columns}

  \takeaway{Collect state, action, and reward for training the local planner.}
\end{frame}

\begin{frame}{Stage 4 | Online learning (Contributions 2 and 3)}
  \source{Zhang et al., IEEE Transactions on Robotics}
  \vspace*{-0.4\baselineskip}
  {\footnotesize\color{body} While walking, train a planner that also uses RGB.}
  \vspace{0.1em}
  \begin{columns}[T,onlytextwidth]
    \begin{column}{0.50\textwidth}
      \infocard[blue][15mm]{Soft Actor Critic}{%
        \li Also use past data
        \li Randomize actions, so one route is not repeated}
      \vspace{0.12em}
      \infocard[orange][15mm]{DroQ}{%
        \li Dropout inside the Q network
        \li Limit overestimation on little data}
      \vspace{0.12em}
      \infocard[green][15mm]{RLPD}{%
        \li Manage the training batch
        \li Half offline data, half online data}
    \end{column}
    \hfill
    \begin{column}{0.50\textwidth}
      \centering
      \includegraphics[width=\linewidth,height=46mm,keepaspectratio]{stage4.png}
    \end{column}
  \end{columns}

  \takeaway{Offline data and randomized actions train efficiently from little data.}
\end{frame}

\begin{frame}{Results | Traversability cost (Contribution 1)}
  \source{Zhang et al., IEEE Transactions on Robotics}
  {\footnotesize\color{body} Replace the training traversability cost $T$ with one of the following.}

  \vspace{0.35em}
  \begin{columns}[T,onlytextwidth]
    \begin{column}{0.32\textwidth}
      \infocard[red][12mm]{Ours}{\li Use the critic value $V$}
    \end{column}
    \hfill
    \begin{column}{0.32\textwidth}
      \infocard[blue][12mm]{Tracking}{\li Velocity-tracking error}
    \end{column}
    \hfill
    \begin{column}{0.32\textwidth}
      \infocard[orange][12mm]{SysID}{\li Friction from a friction estimator}
    \end{column}
  \end{columns}

  \vspace{0.3em}
  \centering
  \includegraphics[width=0.86\textwidth,height=32mm,keepaspectratio]{bars_cost.png}

  \takeaway{Ours walks at the lowest cost.}
\end{frame}

\begin{frame}{Results | Why vision matters (Contribution 2)}
  \source{Zhang et al., IEEE Transactions on Robotics}
  \begin{columns}[T,onlytextwidth]
    \begin{column}{0.32\textwidth}
      \infocard[red][14mm]{Ours}{\li Good on every metric}
    \end{column}
    \hfill
    \begin{column}{0.32\textwidth}
      \infocard[blue][14mm]{RGB only}{\li Many collisions}
    \end{column}
    \begin{column}{0.32\textwidth}
      \infocard[orange][14mm]{Depth only}{\li High traversability cost}
    \end{column}
    \hfill
  \end{columns}

  \vspace{0.35em}
  \centering
  \includegraphics[width=0.78\textwidth,height=34mm,keepaspectratio]{bars_vision.png}

  \takeaway{Depth alone costs more, and RGB alone collides more, so both are needed.}
\end{frame}

\begin{frame}{Results | Online and offline learning (Contribution 3)}
  \source{Zhang et al., IEEE Transactions on Robotics}
  \vspace{0.1em}
  \centering
  {\scriptsize
  \begin{tabular}{lcc}
    \toprule
    & Offline (past examples) & Online (exploration) \\
    \midrule
    \textbf{Ours (RLPD)} & ordinary data 50\% & ordinary data 50\% \\
    Behavior Cloning & imitation data & none \\
    Vanilla SAC & none & ordinary data \\
    SAC with BCL & imitation data & ordinary data \\
    \bottomrule
  \end{tabular}
  }

  \vspace{0.25em}
  \centering
  \includegraphics[width=0.78\textwidth,height=32mm,keepaspectratio]{bars_rlpd.png}

  \takeaway{Ours learns efficiently from little data and reaches high performance.}
\end{frame}

\begin{frame}{Results | Real-robot evaluation}
  \source{Zhang et al., IEEE Transactions on Robotics}
  \centering
  \includegraphics[width=\textwidth,height=54mm,keepaspectratio]{real_cost.png}

  \takeaway{On easy ground, the traversability cost goes down.}
\end{frame}

\begin{frame}{Results | Real-robot navigation}
  \source{Zhang et al., IEEE Transactions on Robotics}
  \centering
  \includegraphics[width=\textwidth,height=54mm,keepaspectratio]{nav_real.png}

  \takeaway{It avoided hard ground and walked to the goal.}
\end{frame}

\begin{frame}{Summary | Traversability-Aware}
  \source{Zhang et al., IEEE Transactions on Robotics}
  \infocard[green][18mm]{Contributions}{%
    \li Traversability from self-experience, with uncertainty
    \li Terrain-avoiding navigation that integrates RGB-D
    \li Online and offline learning on the real robot}
  \vspace{0.35em}
  \infocard[orange][18mm]{Limitations}{%
    \li Adaptation to large height changes, such as stairs}

  \takeaway{Traversability enters the reward. Large height changes are still difficult.}
\end{frame}

\begin{frame}{Summary | The two surveys}
  \cardbalance
  \begin{columns}[T,onlytextwidth]
    \begin{column}{0.49\textwidth}
      \infocard[green][48mm]{Where Should I Walk?}{%
        \li Tie foot contact force to images and predict walking cost
        \li Fill sparse labels with Mean Teacher
        \li Choose the path by Dijkstra on the cost map
        \li Pass walkability to the planner as a map}
    \end{column}
    \hfill
    \begin{column}{0.49\textwidth}
      \infocard[blue][48mm]{Traversability-Aware}{%
        \li Evaluate traversability from self-experience and uncertainty
        \li A local planner that outputs velocity directly from RGB-D
        \li Combine offline and online learning on the real robot
        \li Pass walkability to the action as a reward}
    \end{column}
  \end{columns}
  \takeawaybalance{The first paper puts preference in a map. The second puts it in navigation.}
\end{frame}

\end{document}
